\subsection{CJSR Approximation via SOS}
\label{sec:sos}

The $2d$th root of homogeneous polynomials of degree $2d$ can be used as Lyapunov function.
\begin{proposition}
  \label{prop:pos}
  Consider \defAGe{}.
  Suppose that there exist $|\Nodes|$ strictly positive homogeneous polynomials $p_v(x)$ of degree $2d$
  such that
  \[ p_{v}(A_\sigma x) \leq \gamub^{2d}p_{u}(x) \]
  holds for all \arc{} $(u,v,\sigma) \in \Arcs$.
  Then $\cjsr \leq \gamub$.
  \begin{proof}
    Define $f_v(x) = [p_v(x)]^{\frac{1}{2d}}$ and use \theoref{stabcert}.
  \end{proof}
\end{proposition}

We relax the positivity condition of \propref{pos} by the more tractable sum of squares (SOS) condition and
define $\jsrsos$ as the solution of the following SOS restriction of \progref{primalinf}.
\begin{myprog}[Primal]
  \label{prog:primal}
  \begin{align}
    \notag
    \inf_{p_v(x) \in \Rh[x],\gamub \in \R} \gamub\\
    \notag
    \gamub^{2d} p_u(x) - p_v(A_\sigma x) & \text{ is SOS}, \quad \forall (u, v, \sigma) \in \Arcs,\\
    \label{eq:soscons1}
    p_v(x) & \text{ is SOS}, \quad \forall v \in \Nodes,\\
    \label{eq:soscons3}
    p_v(x) & \text{ is strictly positive}, \quad \forall v \in \Nodes,\\
    \notag
    \sum_{v \in V} \int_{\mathbb{S}^{n-1}} p_v(x) \dif x & = 1.
  \end{align}
\end{myprog}

\begin{myrem}
  In practice we can replace \eqref{eq:soscons1} and \eqref{eq:soscons3} by ``$p_v(x) - \epsilon\|x\|_2^{2d}$ is SOS''
  for any $\epsilon > 0$.
  This constrains $p_v(x)$ to be in the interior of the SOS cone,
  which is sufficient for $p_v(x)$ to be strictly positive.
  The bounds given in \Sect\ref{sec:quality} are valid
  if $p_v(x)$ is in the interior of the SOS cone.
\end{myrem}

By \propref{pos},
a feasible solution of Program~\ref{prog:primal} gives an upper bound for $\cjsr$,
and thus, for any positive degree $2d$,
\begin{equation}
  \label{eq:upperbound}
  \cjsr \leq \jsrsos.
\end{equation}


\begin{myexem}
  \label{exem:simple1}
  Consider the unconstrained system \cite[Example~2.1]{ahmadi2012joint} with $m=3$:
  \[ \A = \{ A_1 = e_1e_2^\Tr , A_2 = e_2e_3^\Tr , A_3 = e_3e_1^\Tr  \} \]
  where $e_i$ denotes the $i$th canonical basis vector.

  One can verify that for any $d$, a solution to Program~\ref{prog:primal} is given by
  \[ (p(x), \gamma) = (x_1^{2d} + x_2^{2d} + x_3^{2d}, 1). \]
\end{myexem}

\begin{myexem}
  \label{exem:run2}
  Let us reconsider our running example; see Example~\ref{exem:run1}.
  The optimal solution of Program~\ref{prog:primal}
  is represented by
  Figure~\ref{fig:primal} for $2d = 2$, 4, 6, 8, 10 and 12.
  \begin{figure}[!ht]
    \centering
    \tabulinesep=0.5mm
    \begin{tabu}{{X[0.85,c,m]X[c,m]X[c,m]X[c,m]X[c,m]}}
      & Node 1 & Node 2 & Node 3 & Node 4\\
      \PrimalFig{2}{1}
      \PrimalFig{4}{2}
      \PrimalFig{6}{3}
      \PrimalFig{8}{4}
      \PrimalFig{10}{5}
      \PrimalFig{12}{6}
    \end{tabu}
    \caption{
    Representation of the solutions to Program~\ref{prog:primal}
    with different values of $d$ for the running example.
    The blue curve represents the boundary of the
    1-sublevel set of the optimal
    solution $p_v$ at each node $v \in \Nodes$.
    The dashed curve is the boundary of the unit circle.
    Observe that some sets are not convex.}
    \label{fig:primal}
  \end{figure}
\end{myexem}

